\section{磁场与磁感强度}
在本章开始时，我们用了相当大的篇幅来强调，磁是运动电荷在电之外的那一部分效应，但是这样的磁似乎与我们熟悉的磁铁的那种磁并不相符，这实际上是一个历史问题。 

在历史上，有关磁现象的研究要远早于电现象，在公元前6世纪我国古代先民就已经懂得用磁石吸引铁器，在公元18世纪后西方科学家才开始逐渐研究电学现象。而磁现象的发现正是始于磁铁的磁性，我们观察到无论是天然磁铁还是人造磁铁，都具有吸引铁、镍、钴及其合金的性质，这种性质称为\uwave{磁性}，这些能被磁铁吸引的物质则统称为\uwave{铁磁质}。

我们发现，条形磁铁两端磁性特别强，称为\uwave{磁极}，条形磁铁中间磁性很弱，称为\uwave{中性区}，另外一个有趣的现象是，如果我们将条形磁铁悬挂起来，使它能在水平面内自由转动，我们观察到，无论我们条形磁铁的最初指向是什么，在稳定后两端的磁极总是分别指向南北方向，并且，总是一个磁极指向南方，而另外一个磁极指向北方，这就说明两个磁极是有差异的。

为了便于区分磁铁的磁极，将磁铁指北的一端称为\uwave{北极}，将磁铁指南的一端称为\uwave{南极}，两者分别可以用字母N和S表示，基于这样的磁极定义，我们还发现如果我们尝试将两根磁铁相互靠近，我们将会看到这样的现象，\textbf{同号磁极相互排斥，异号磁极相互吸引}，这一实验现象说明，地球本身就是一个巨大的磁铁，北极是它的磁南极，南极是它的磁北极。不过实际上地球的地理南北极与磁南北极并不是完全重合的，两者间的夹角称为\uwave{地磁偏角}。

在1785年库仑提出关于\uwave{电的库仑定律}，即我们熟悉的\Refx{定律}{库仑定律}的同时，也提出了\uwave{磁的库仑定律}，它认为磁现象是由于\uwave{磁荷}产生的，就像电现象是由于电荷产生的
\begin{itemize}
    \item 电荷将会产生电场，电荷有“正负”两种，同号电荷相斥，异号电荷相吸，电荷间的相互作用力，与两者电荷量的乘积成正比，与电荷间距离的平方成反比。
    \item 磁荷将会产生磁场，磁荷有“北南”两种，同号磁荷相斥，异号磁荷相吸，磁荷间的相互作用力，与两者磁荷量的乘积成正比，与磁荷间距离的平方成反比。
\end{itemize}
在磁荷磁的库仑定律的帮助下，有关磁铁的现象基本可以被定性解释和定量计算。

\begin{Figure}{磁的磁荷观点}
    \subfigure[磁化前]
    {
        \inputfigure[scale=0.45]{Chapter10A_Figure03}
    }\\
    \subfigure[磁化后]
    {
        \inputfigure[scale=0.45]{Chapter10A_Figure02}
    }\\
    \subfigure[磁化的宏观效果]
    {
        \inputfigure[scale=0.45]{Chapter10A_Figure04}
    }
\end{Figure}

磁荷理论中，磁铁当然不是由两个巨大的“北磁荷”和“南磁荷”拼接而成的，且不深究两者是如何克服巨大斥力结合在一起的，我们知道，磁铁的一个非常奇特的性质就是，如果我们将磁铁沿中性区切成两块，那么我们不会得到两个独立的南北极，而是产生了两个具有完整南北极的小磁铁。因此实际上，磁铁在磁荷理论中被认为是这样一种物质，其内部存在有大量排列整齐的微小磁铁，这种微小磁铁也称为\textbf{磁偶极子}，这类似于电场中的\textbf{电偶极子}。

磁铁的中间部分，磁偶极子的南极和北极总是相互抵消，从而对外在宏观上不显出磁性，这就是所谓的中性区，磁铁两端处位于端面附近的那些磁偶极子，具有大量无法抵消的北磁荷和南磁荷，从而就形成了具有强烈的磁性的北磁极和南磁极，就像被极化的电介质那样。

磁铁的任何一个磁极都可以吸引铁磁质，而铁磁质本身没有任何磁性。这就是因为铁磁质内部同样存在大量的磁偶极子，但是这些磁偶极子原先排列非常混乱，从而在宏观上不显示出任何磁性。而在磁铁的外加磁场作用下，铁磁质中的磁偶极子排列变得整齐，铁磁质就被磁化了，这个过程与电介质中的电偶极子在电极化的过程中排列整齐是十分相似的。但是这里的铁磁质作为一种磁介质按照原先电介质的观点看是较为特殊的，因为在磁化后，即便外磁场褪去，这些\textbf{铁磁质}也往往可以保留一定磁性，这相当于电介质中的\textbf{铁电体}。有一部分铁磁质时间长了或是经历碰撞和加热之后就会被消磁，有一部分铁磁质则可以更为稳定的长久保留磁性，这类磁介质也称为\textbf{永磁体}，这相当于电介质中的\textbf{永驻体}，实际上，我们所使用的天然磁铁就是一种天然的永磁体，而所谓的人造磁铁就算通过磁化的方式制造的永磁体。

电学曾定义了电场强度$\vb*{E}$，磁学中也可以相仿的定义磁场强度$\vb*{H}$，即单位磁荷在磁场中所受到的力（我们或许需要将北磁荷和南磁荷分别赋予正和负），通过点磁荷的库仑定律可以先推导出点磁荷的磁场，进而计算出磁偶极子的磁场，这样就可以定量分析磁铁了。

\textbf{然而，如果以上这一切都是真的，磁学就会比今天所看到的简单太多了。}

磁荷理论的一个缺陷在于，\uwave{磁单极子}从未被发现，即独立的北磁荷和南磁荷并不存在，如果我们将一个磁铁切开只会得到更多的小磁铁，也没有任何一种粒子具有单独的磁荷。而相比之下在电学中，\uwave{电单极子}则是常见的，即独立的正电荷和负电荷，质子和电子就是最容易看到的例子。然而实际上，磁单极子不存在并非是磁荷理论被历史舍弃的主要原因，因为我们完全可以退一步认为所有磁荷都是以类似于束缚电荷那样的束缚磁荷的形态存在的。

磁荷理论的真正问题在于，电流和运动电荷的磁效应的发现。

在1820年，丹麦科学家奥斯特通过著名的\uwave{奥斯特实验}，发现将一个可以自由旋转的小磁针置于直导线的正下方，然后给直导线通以电流，观察到小磁针将发生偏转，而且当电流方向改变时，观察到小磁针的偏转方向将相应的发生反转，\textbf{这就说明了电流同样可以产生磁场}。

在这里有一个非常相似的实验，如果将一个根水平导悬挂在一个马蹄形磁铁的磁极之间，然后给直导线通以电流，观察到水平导线将会发生偏移，\textbf{这就说明了电流会受到磁场作用}。

在这个跨时代的实验的消息抵达法国科学院的仅1周之后，法国科学家安培通过实验，发现两根平行载流导线之间也存在相互作用力，如果导线的电流方向相同则相互吸引，如果导线的电流方向相反则相互排斥，\textbf{这就说明电流之间可以独立的通过磁场产生相互作用}。

电流是电荷定向运动形成的，通过实验我们也可以观察到磁场对运动电荷的作用，阴极射线管是一个真空放电管，在它两个电极之间加上高电压时，在它的阴极就会产生所谓的阴极射线，也就是电子束。电子束本身是不可见的，为此在阴极射线管中附有可以观测电子束轨迹的荧光屏，通常情况下电子束是沿直线前进的，但是如果我们在阴极射线管旁放置一块磁铁，就会观察到电子束的轨迹发生了偏转，这就可以表明电子束的受到了磁场作用。

磁铁和磁铁间，电流和磁铁间，电流和电流间，都存在相互作用，当然或许这其实是三种不同的作用，但是从自然规律的统一性上看我们更愿意相信这实际是同一种磁相互作用，然而这就导致了一个问题，尽管我们可以称磁铁中有一些磁偶极子，但是电流中绝无可能有任何的磁偶极子，也就是说，磁荷理论无法解释电流的磁效应，至此，磁荷理论也就被抛弃了。

在今天，我们从一个相反的观点来看待磁现象，我们认为磁场的本源是电流或者说是电流中的运动电荷，这也就是本章开始处“磁是运动电荷在电之外的那一部分效应”的含义，那么在这种新的认识下，磁铁的磁效应又应该如何解释呢？安培提出了分子电流假说，这一假说认为，磁铁在分子层面中有许多排列整齐的微小环状电流，这些取向一致分子电流在宏观上等效于一个环绕磁铁边缘的大环状电流，从而对外显示出磁性，磁铁的磁化实际上就是将原先铁磁质中排列混乱的分子电流通过磁场作用排列整齐，从而使得铁磁质也具有磁性。
\begin{Figure}{磁的分子电流观点}
    \subfigure[磁化前]
    {
        \inputfigure[scale=0.45]{Chapter10A_Figure06}
    }\\
    \subfigure[磁化后]
    {
        \inputfigure[scale=0.45]{Chapter10A_Figure05}
    }\\
    \subfigure[磁化的宏观效果]
    {
        \inputfigure[scale=0.45]{Chapter10A_Figure07}
    }
\end{Figure}

安培分子电流假说与实验事实是相符的，因为依照它的观点，磁铁像是一个由分子电流构成的通电螺旋管，而实际上，磁铁产生的磁场与真正的通电螺旋管的磁场是非常相似的。

安培分子电流假说不仅与电流的磁效应兼容，而且其与实际的物质微观结构是相符的，我们现在知道，电子在环绕原子核转动的同时，电子本身也会发生自旋，这两者实际就是分子电流的微观来源。当然这种“电子绕核转动”和“电子自旋”只是经典观点下提法，现在我们已经知道，电子实际是以电子云的形态存在，电子自旋也只是一个粒子属性，并非是像真正的宏观物体那样的自转，这里最为本质的研究需要用到量子力学，但是经典力学的这些提法具有启发性的，尽管在原理上它不完全正确，但足矣对物质磁性给出足够准确的猜测了。

本章将首先研究真空中的磁场，而有关那些磁介质和磁铁以及安培分子电流假说的内容将在下一章讨论。暂时我们只需要理解，即便现在我们认为磁场是源于电流，磁铁的磁场仍然是同样客观真实的，磁铁的磁效应最终也可以用电流解释，并非什么理论之外的奇怪东西。

\subsection{磁感应强度}
在电场中，我们引入了描述静止电荷间的电相互作用力的\Refx{定律}{库仑定律}，随后通过该结论推导出试探电荷在空间中某一特定位置受到的电场力仅与其电荷量有关，即
\begin{Equation}[电场力的作用]
    \vb*{F}=q\vb*{E}
\end{Equation}
在磁场中，由于一些原因，我们很难通过实验测定运动电荷间的磁相互作用力的公式，但是我们发现，如果现在有一些磁场，或许是由旁边的电流和磁铁产生的，这不重要，那么通过实验将观察到试探电荷在空间某一特定位置受到的磁场力仅与其电荷量和速度有关，即\footnote{这里要假定空间中没有任何电场，或者我们能计算出电场力的大小并在试探电荷的受力中减去，因为实验中我们并不能分辨试探电荷受到的作用力中那一部分属于电场力那一部分属于磁场力。这里研究的是磁场，我们只需要其中磁场力的那一部分。}
\begin{Equation}[磁场的作用]
    \vb*{F}=q\vb*{v}\times\vb*{B}
\end{Equation}
式\ref{式:磁场的作用}要比式\ref{式:电场力的作用}复杂，电场力仅和电荷量有关，磁场力同时与电荷量和速度有关。

式\ref{式:磁场的作用}中的矢量$\vb*{B}$与先前电场中的电场强度$\vb*{E}$的地位是相似的，其定量反映了磁场的空间分布，称为\uwave{磁感应强度}或简称为\uwave{磁感强度}。磁感强度的确定可以通过测量试探电荷在磁场中的磁场力实现，令试探电荷$q$以某个大小相同的速度$\vb*{v}$经过同一位置，将试探电荷$q$受到的最大磁场力记为$F_{\max}$，根据叉乘原理，此时$\vb*{F}_{\max},\vb*{v},\vb*{B}$垂直且满足$F_{\max}=qvB$，故
\begin{BoxDefinition}[磁感强度]
    \vb*{B}=\frac{\vb*{F}_{\max}\times\vu*{v}}{qv}
\end{BoxDefinition}
这里的矢量关系可以由下图看出，即当磁场力最大$\vb*{F}=\vb*{F}_{\max}$时，此时磁感强度$\vb*{B}$垂直于由$\vb*{F},\vb*{v}$构成的平面，且磁感应强度$\vb*{B}$的方向为右手四指由$\vb*{F}$转向$\vb*{v}$时大拇指的方向。
\begin{Figure}{磁感强度的矢量关系}
    \subfigure[当$\vb*{v},\vb*{B}$为垂直时的$\vb*{F}$视角]
    {
        \hspace{30pt}
        \inputfigure{Chapter10B_Figure01}
        \hspace{30pt}
    }
    \subfigure[当$\vb*{v},\vb*{B}$为垂直时的$\vb*{B}$视角]
    {
        \hspace{30pt}
        \inputfigure{Chapter10B_Figure02}
        \hspace{30pt}
    }\\
    \subfigure[当$\vb*{v},\vb*{B}$不垂直时的$\vb*{F}$视角]
    {
        \hspace{30pt}
        \inputfigure{Chapter10B_Figure03}
        \hspace{30pt}
    }
    \subfigure[当$\vb*{v},\vb*{B}$不垂直时的$\vb*{B}$视角]
    {
        \hspace{30pt}
        \inputfigure{Chapter10B_Figure04}
        \hspace{30pt}
    }
\end{Figure}

这里我们也可以看到，在一般情况下，实际上$\vb*{B}$并不总是垂直于$\vb*{F},\vb*{v}$构成的平面，因为先验成立的是$\vb*{F}=q\vb*{v}\times\vb*{B}$，所以只有当三者垂直磁场力最大时我们才能逆向求出磁感强度。

磁感强度在国际单位制中的单位是\si{T=N\cdot A^{-1}\cdot m^{-1}}，称为\uwave{特斯拉}，量纲是\si{[MT^{-2}I^{-1}]}。

\textbf{磁感强度}描述了磁场分布，虽然其地位等同于电场强度，但是不能将其称为\textbf{磁场强度}，这就是因为本节初叙述的历史原因，磁场强度$\vb*{H}$的名称已经在磁荷理论中被先行被定义为了单位磁荷在磁场中受到的力，因而后来基于试探运动电荷的定义只能改称为磁感强度$\vb*{B}$。

\subsection{磁感强度的叠加原理}
实验表明，当空间中有若干个磁场源的情况下，它们产生的磁场服从叠加原理，它们的总磁感强度就是各个磁场源独立存在时的磁感强度$\vb*{B}_1,\vb*{B}_2,\cdots,\vb*{B}_n$的矢量和，即
\begin{BoxPrinciple}[磁感强度的叠加原理]
    若干磁场源在磁场中某点处产生的磁感强度，

    等个各磁场源单独存在时，在该点处产生的磁感强度的矢量和。
    \begin{BoxEq}
        \vb*{B}=\SumIN{\vb*{B}_i}
    \end{BoxEq}
\end{BoxPrinciple}
\Refx{定律}{磁感强度的叠加原理}与先前\Refx{定律}{电场强度的叠加原理}是电磁场中最基本的两个规律，电场中的这一点是我们从力的叠加原理推证出来的，磁场中如果我们愿意的话其实仍然可以这样做，但实际上，场的叠加原理才是更为普遍和基本的规律。\footnote{场的叠加原理是更普遍和基本的，因此从力的叠加原理推导它实在没有什么必要，唯一的价值是说明两者不矛盾。}

\section{毕奥-萨伐尔定律}

\subsection{毕奥-萨伐尔定律}
本节我们来研究运动电荷或运动电荷构成的电流是如何产生磁场的，然而由于一些后面我们会提到的原因，运动电荷产生的磁场很难被实验观测到，运动电荷构成的电流又往往是具有特定形状的，例如通电直导线或通电线圈，无法归纳出普适性比较强的规律。

在安培发现两根载流导线之间的磁相互作用力后，另外两位法国科学家毕奥和萨伐尔随后提出了定量描述电流磁场的\uwave{毕奥-萨伐尔定律}。这里引入了\uwave{电流元}的概念，我们知道，由于磁场适用叠加原理，我们就可以将任意形状的电流视为无穷多段无限小的电流的集合，这些无限小的电流就是所谓的电流元，因此只要能弄清电流元产生的磁场（这就是毕奥-萨伐尔定律的具体内容），那么任意形状的电流的磁场就可以通过电流元的磁场的叠加求得了。

我们需要指出的是，电流元虽然长度无限小，但是仍然是有方向的，故通常记作$I\dif\vb*{l}$。

我们现在来看毕奥-萨伐尔定律是如何表述的
\begin{BoxGeneral}[毕奥-萨伐尔定律][定律]
    \begin{BoxEq}
        \dif{\vb*{B}}=\frac{\mu_0}{4\pi}
        \frac{I\dif{\vb*{l}}\times\vu*{r}}{r^2}
    \end{BoxEq}
\end{BoxGeneral}
\Refx{定律}{毕奥-萨伐尔定律}指出，电流元产生的磁场垂直于$I\dif{\vb*{l}},\vb*{r}$构成的平面，且指向右手四指由$I\dif{\vb*{l}}$转向$\vb*{r}$时大拇指的方向，其大小，与距离平方成反比，与电流成正比。

\Refx{定律}{毕奥-萨伐尔定律}中$\mu_0$为真空磁导率，其数值为$\mu_0=4\pi\times 10^{-7}$\si{N\cdot A^{-2}}，这里真空磁导率的数值之所以如此齐整，是因为和安培的定义有关，稍后我们会进一步讨论。

\Refx{定律}{毕奥-萨伐尔定律}并不总是成立的，其仅在\uwave{静磁场}的情况下精确成立，这实际上要求产生磁场的电流是一个恒定电流场，即组成电流的每个电流元的大小和方向不随时间变化，这与先前电荷间的\Refx{定律}{库仑定律}只能适用于\uwave{静电场}的情况是类似的。

\Refx{定律}{毕奥-萨伐尔定律}可以根据磁场叠加原理计算任意载流导线的磁场
\begin{BoxEquation}[电流的磁场]
    \vb*{B}=\Ilnt[L]{\dif\vb*{B}}=\Ilnt[L]{\frac{\mu_0}{4\pi}
    \frac{I\dif{\vb*{l}}\times\vu*{r}}{r^2}}
\end{BoxEquation}
这里只需要曲线积分，因为目前我们考虑的电流都是在导线中的，即电流都是线状的。

值得注意的是，\Refx{定律}{毕奥-萨伐尔定律}中设想的电流元只是一种微元模型，电流元是不可能单独存在的，因此该定律是无法直接通过物理实验验证的，但是通过其积分得到的各类不同载流导线的磁场都与实验都相符，这就间接验证了定律的正确性。

\subsubsection{载流直导线的磁场}
设有一载流直导线，长度为$L$，电流为$I$，现考察载流直导线外的场点$P$处的磁感强度。

设$P$点到导线的垂足$O$为原点，以直导线上的电流方向为$x$轴的正方向建立平面直角坐标系，在该坐标系中点$P$在$y$轴上，且设其离开导线的垂直距离（即纵坐标）为$y$。

设$P$点和导线两端的连线与导线之间的夹角分别为$\theta_1$和$\theta_2$。
\begin{Figure}{载流直导线的磁场}
    \inputfigure[scale=0.85]{Chapter10B_Figure05}
\end{Figure}
我们可以将载流直导线视为一系列的电流元的组合，根据\Refx{定律}{毕奥-萨伐尔定律}可知，导线上方的磁场垂直直面向外，导线下方的磁场垂直直面向内，即在同一场点处，每个电流元产生的磁场方向都相同，因此我们就可以将矢量$\dif{\vb*{B}}$转化为标量$\dif{B}$计算。

这里为了方便在平面图形上表示那些垂直于纸面的方向，如\Refx{图}{载流直导线的磁场}所示，我们往往会用$\odot$和$\otimes$表示\uwave{垂直纸面向外}和\uwave{垂直纸面向内}的概念\footnote{实际上这种记号在磁场更多是用于表达垂直直面的电流，而不是磁感强度。}，这种记号可以通过离弦的箭来记忆，射向自己的箭将看到点状箭尖，射出的箭将看到十字形的尾部箭羽。

这里以载流直导线上方的场点$P$为例，将\Refx{定律}{毕奥-萨伐尔定律}改写为标量形式
\begin{Equation}[载流直导线1]
    \dif{B}=\frac{\mu_0}{4\pi}\frac{I\dif{l}\sin\theta}{r^2}
\end{Equation}
式\ref{式:载流直导线1}无法直接积分，因为其关于$\theta$但却具有微分$\dif{l}$，因此要首先将$\dif{l}$转换为$\dif{\theta}$。

关注到$l=r|\cos\theta|$和$y=r|\sin\theta|$，因此
\begin{equation*}
    \frac{l^2}{y^2}=\cot^2\theta\qquad\frac{l}{y}=\abs{\cot\theta}
\end{equation*}
因此$r^2$可以被表示为
\begin{Equation}[载流直导线2]
    r^2=y^2+l^2=y^2(1+\cot^2\theta)=y^2\csc^2\theta
\end{Equation}
因此$l$的微分$\dif{l}$可以被表示为
\begin{Equation}[载流直导线3]
    l=y|\cot\theta|\qquad \dif{l}=y|\csc^2\theta\dif{\theta}|=y\csc^2\theta\dif{\theta}
\end{Equation}
这就有
\begin{Equation}[载流直导线4]
    \frac{I\dif{l}}{r^2}=\frac{Iy\csc^2{\theta}}{y^2\csc^2\theta}=\frac{I\dif{\theta}}{y}
\end{Equation}
式\ref{式:载流直导线4}代入式\ref{式:载流直导线1}中，得到
\begin{Equation}[载流直导线5]
    \dif{B}=\frac{\mu_0}{4\pi}\frac{I\sin{\theta}\dif{\theta}}{y}
\end{Equation}
这就将自变量$\theta$和微分$\dif{\theta}$相统一。

对$\dif{B}$积分，得到
\begin{equation*}
    B=\Int[\theta_1][\theta_2]{\dif{B}}=
    \frac{\mu_0}{4\pi}\frac{I}{y}
    \Int[\theta_1][\theta_2]{\sin{\theta}\dif{\theta}}=
    \frac{\mu_0}{4\pi}\frac{I}{y}(\cos\theta_2-\cos\theta_1)
\end{equation*}
即
\begin{BoxEquation}[载流直导线的磁场]
    B=\frac{\mu_0}{4\pi}\frac{I}{y}(\cos\theta_2-\cos\theta_1)
\end{BoxEquation}
这里，如果载流直导线变为无限长，那么$\theta_1\rightarrow 0,\theta_2\rightarrow\pi$
\begin{BoxEquation}[载流无限长直导线的磁场]
    B=\frac{\mu_0}{2\pi}\frac{I}{y}
\end{BoxEquation}
这里我们可以通过\Refx{图}{载流直导线的磁场}将载流直导线附近的磁场方向归纳如下，即将右手大拇指指向电流方向，那么右手四指的环绕方向即为磁场方向，称为\uwave{右手螺旋法则}。

\subsubsection{载流圆线圈的磁场}
设有一载流圆线圈，半径为$R$，电流为$I$，现考察其轴线上的场点$P$处的磁感强度。
\begin{Figure}{载流圆线圈的磁场}
    \inputfigure[scale=0.64]{Chapter10B_Figure06}
\end{Figure}
设线圈的圆心$O$为原点，以轴线为$x$轴，以线圈所在平面为$yz$平面，建立空间直角坐标系，场点$P$至$yz$平面的距离为$x$，场点$P$至线圈任意一点的距离均为$r=\sqrt{R^2+x^2}$。

这里$\vb*{I}\dif{\vb*{l}},\vb*{r}$总是垂直的，因而根据\Refx{定律}{毕奥-萨伐尔定律}
\begin{equation*}
    \dif{B}=\mk\frac{I\dif{l}}{r^2}
\end{equation*}

如\Refx{图}{载流圆线圈的磁场}所示，由于线圈的对称性，在线圈各处的电流元对轴线上场点施加的磁感强度中，垂直$x$轴的各分量相互抵消，只有平行$x$轴的分量才有贡献。

由于$\cos\alpha=\dfrac{R}{r}$，且$r=\sqrt{R^2+x^2}$
\begin{equation*}
    \dif{B}_x=\dif{B}\cos\alpha=\mk\frac{I\dif{l}}{r^2}\frac{R}{r}=\mk\frac{IR\dif{l}}{r^3}=\mk\frac{IR\dif{l}}{(R^2+x^2)^{\frac{3}{2}}}
\end{equation*}
从而积分得
\begin{equation*}
    B=B_x=\mk\frac{IR\dif{l}}{\qty(R^2+x^2)^{\frac{3}{2}}}\Int[0][2\pi R]{\dif{l}}=
    \frac{\mu_0}{2}\frac{IR^2}{\qty(R^2+x^2)^{\frac{3}{2}}}
\end{equation*}
即
\begin{BoxEquation}[载流圆线圈的磁场]
    B=\frac{\mu_0}{2}\frac{IR^2}{\qty(R^2+x^2)^{\frac{3}{2}}}
\end{BoxEquation}
对于场点位于线圈圆心的情况，这里有$x=0$，即
\begin{BoxEquation}[载流圆线圈在圆心处的磁场]
    B=\frac{\mu_0}{2}\frac{I}{R}
\end{BoxEquation}
对于场点远离线圈的情况，这时有$x\gg R$，即
\begin{BoxEquation}[载流圆线圈在远场点sz的磁场]
    B=\frac{\mu_0}{2}\frac{IR^2}{x^3}
\end{BoxEquation}
以上我们是以场点$P$位于线圈上方的情况考虑的，但是通过\Refx{图}{载流圆线圈的磁场}的下半部分可以看出，这样的推导对线圈下方的场点也是成立的，而且无论场点在线圈上方还是线圈下方，其磁场方向都是沿着线圈电流的右手螺旋方向，这就为矢量表达创造了机会。

记线圈的面积为$S=\pi r^2$，定义载流线圈的\uwave{磁偶极矩}为
\begin{BoxDefinition}[磁偶极矩]
    \vb*{m}=IS\vb*{e_n}
\end{BoxDefinition}
其中$\vb*{e}_n$为载流线圈电流的右手螺旋方向上的单位矢量。
\begin{Figure}{载流圆线圈与磁偶极子}
    \inputfigure[scale=0.64]{Chapter10B_Figure07}
\end{Figure}

那么对于载流线圈的远场点，磁感强度就可以表示为
\begin{BoxEquation}[磁偶极子在中轴线上的磁场]
    \vb*{B}=\frac{\mu_0}{2\pi}\frac{\vb*{m}}{x^3}
\end{BoxEquation}
\Refx{公式}{磁偶极子在中轴线上的磁场}与先前\Refx{公式}{电偶极子延长线上的电场}在形式上很接近，出于这样的原因，仿照电偶极子，我们也将载流线圈称为\uwave{磁偶极子}，因此我们在这里的\Refx{定义}{磁偶极矩}在地位上就与先前的\Refx{定义}{电偶极矩}是相当的了。

磁偶极子和电偶极子在两者偶极矩方向上的远场点处的场的数学形式非常相似，这一点已经通过数学推导证明了，但在下一节我们还会通过场线图的对比看到，磁偶极子和电偶极子在任意远场点处的场的形态也都非常相似。我们或许会感到困惑，因为虽然电场和磁场都是平方反比场，但是两者的性质完全不同，电场是有源无旋场，磁场是有旋无源场，为何两种性质截然不同的场在这里可以产生相同的数学形式？这实际就和偶极子中远场点的要求是密不可分的，因为当距离很远时，无论是源还是旋都会归于零，从而消除了差异。

磁偶极子可以被称为“偶极子”的另一个原因是，载流线圈与先前磁荷理论中同样称为磁偶极子的小磁铁产生的磁场是相当的，即，\textbf{载流线圈的磁场与小磁铁的磁场是等效的}，这再次体现了磁荷理论与分子电流理论在解释磁铁的现象时的一致性，磁铁的磁性总是来自于材料中微小的磁偶极子，无论你认为磁偶极子到底是微小的载流线圈还是微小的磁铁。\footnote{载流线圈的磁场反映的其实就是环状电流的磁场，只不过微观上的分子环状电流其实已经不再需要线圈作为载体了。}
\begin{Figure}{载流线圈与小磁铁的磁场对比}
    \subfigure[载流线圈的磁场]
    {
        \includegraphics[scale=0.8]{old/Python/build/VFP_B03.pdf}
    }\qquad
    \subfigure[小磁铁的磁场]
    {
        \includegraphics[scale=0.8]{old/Python/build/VFP_M03.pdf}
    }
\end{Figure}
\vspace{-20pt}

\subsubsection{载流螺线管的磁场}
螺线管是指均匀地密绕在直圆柱面上的螺旋形线圈，设螺线管的半径为$R$，设螺线管中通有电流$I$，沿轴向单位长度的线圈数为$n$，现在来考察其管内轴线上一点的磁场。

螺线管在一般情况下的磁场是较为复杂的，为了便于计算，我们做几点假定。

我们不难想象，对于真实的螺线管，导线每环绕圆柱面一周，导线在轴线方向就要前进一些距离，这在数学上处理起来就很复杂了，但是如果螺线管上的导线缠绕的非常紧密，导线在同一周上的轴线位移可以忽略，这就使得我们可以将完整的螺旋线视为由一系列分离的圆线圈构成，这就是为何我们将螺线管的剖视图绘制为\Refx{图}{载流螺线管的磁场}的原因。

我们还要假定的一点时，螺旋管上的导线缠绕的如此之紧密，以至于即便我们沿着轴线方向取很小的一段距离，在其上的线圈数目也是十分多的。在这里我们实质是将螺线管上的分立线圈数目通过物理无穷小的想法连续化了，目的是为了使得我们可以通过积分方法叠加各个线圈产生的磁场。当然，这会忽略在导线直径尺度上磁场的不连续性，不过现在的螺线管的线径往往是很细的，在几厘米的距离上就可以缠绕数千圈，这种误差便也可以接受了。

\begin{Figure}{载流螺线管的磁场}
    \inputfigure[scale=0.8]{Chapter10B_Figure08}\vspace{-15pt}
\end{Figure}

以$P$为原点，以螺线管电流的右手螺旋方向为$x$轴，建立平面直角坐标系。现在我们在螺线管上距离$P$点为$l$处的位置取一小段$\dif{l}$上的线圈，在这$\dif{l}$上共有线圈$n\dif{l}$匝，那么根据\Refx{公式}{载流圆线圈的磁场}，在这$\dif{l}$上的线圈在场点$P$处产生的磁感强度$\dif{B}$为
\begin{equation*}
    \dif{B}=\frac{\mu_0}{2}\frac{InR^2\dif{l}}{\qty(R^2+l^2)^{\frac{3}{2}}}
\end{equation*}

在螺线管上各个小段$\dif{l}$上的线圈在场点$P$处产生的磁感强度$\dif{B}$方向均相同，指向电流的右手螺旋方向，因此整个螺线管的磁场即为（设$l_1,l_2$是以$P$为中心螺线管两端的横坐标）
\begin{Equation}[载流螺线管1]
    B=\Int[l_1][l_2]{\frac{\mu_0}{2}\frac{InR^2\dif{l}}{\qty(R^2+l^2)^{\frac{3}{2}}}}
\end{Equation}
式\ref{式:载流螺线管1}关于$\dif{l}$的积分较难积出，为此我们使用三角换元，假设位于$l$处的螺线管相对于场点$P$处的角度为$\beta$，根据螺线管中的几何关系可以看到
\begin{equation*}
    l=R\cot\beta\qquad \dif{l}=-R\csc^2\beta\dif{\beta}
\end{equation*}

另外一方面
\begin{equation*}
    l^2+R^2=R^2\csc^2\beta
\end{equation*}

这样就有
\begin{Equation}[载流螺线管2]
    \frac{InR^2\dif{l}}{\qty(R^2+l^2)^{\frac{3}{2}}}
    =-\frac{InR^3\csc^2\beta\dif{\beta}}{R^3\csc^3\beta}
    =-\frac{In\dif{\beta}}{\csc\beta}
    =-In\sin\beta\dif{\beta}
\end{Equation}
假设$l_1,l_2$对应的角为$\beta_1,\beta_2$，那么
\begin{Equation}[载流螺线管3]
    B=\frac{\mu_0}{2}In\Int[\beta_1][\beta_2]
    {-\sin{\beta}\dif{\beta}}
\end{Equation}
即
\begin{BoxEquation}[载流螺线管的磁场]
    B=\frac{\mu_0}{2}In(\cos\beta_2-\cos\beta_1)
\end{BoxEquation}
当螺线管的长度远远大于螺线管的半径，即$L\gg R$时，此时得到了一个无限长螺线管。

对于无限长螺旋管的内部各点处，有$\beta_1\rightarrow\pi$而$\beta_2\rightarrow 0$
\begin{BoxEquation}[载流无限长螺线管的轴线磁场]
    B=\mu_0nI
\end{BoxEquation}

\Refx{公式}{载流无限长螺线管的轴线磁场}指出，无限长载流螺线管轴线上的磁感强度与位置无关，也就是说，载流无限长螺线管在轴线上将产生匀强磁场，这是很重要的结论。\footnote{我们后面还会进一步证明，实际上无限长螺线管内部的磁场都是匀强磁场，不仅限于螺线管的轴线。}

对于半无限长螺线管的端点处，有$\beta_1=\pi/2$和$\beta_2\rightarrow 0$
\begin{BoxEquation}[载流半无限长螺线管的端点磁场]
    B=\frac{1}{2}\mu_0nI
\end{BoxEquation}
因而对于有限长的螺线管，其中央的磁感强度为$\mu_0nI$，其端点的磁感强度为$\mu_0nI/2$，即螺线管中央的磁感强度为螺线管端点处的磁感强度的两倍，且螺线管中间部分可以视为匀强磁场，当然以上这都是需要近似的结果，显然当$L$与$R$比值越大时近似程度越好。

我们先前在\Refx{公式}{载流螺线管的磁场}的推导中，由于三角换元的原因我们将结论转化为了基于$\beta_1,\beta_2$来表述，但很明显更为直观的办法通过场点的横坐标$x$表达。

我们为此将\Refx{图}{载流螺线管的磁场}的原点固定至螺线管中心。
\begin{Figure}{载流螺线管的中心坐标}
    \inputfigure[scale=0.8]{Chapter10B_Figure10}\vspace{-15pt}
\end{Figure}\vspace{-10pt}
由上图可以看出$x,\beta_1,\beta_2$的关系
\begin{Equation}
    \cos\beta_1=
    -\frac{x+L/2}{\sqrt{R^2+(x+L/2)^2}}\qquad
    \cos\beta_2=
    -\frac{x-L/2}{\sqrt{R^2+(x-L/2)^2}}
\end{Equation}
将上式代入\Refx{公式}{载流螺线管的磁场}
\begin{BoxEquation}[载流螺线管的磁场的坐标表示]
    B=\frac{\mu_0}{2}In\qty(\frac{x+L/2}{\sqrt{R^2+(x+L/2)^2}}-\frac{x-L/2}{\sqrt{R^2+(x-L/2)^2}})
\end{BoxEquation}
将\Refx{公式}{载流螺线管的磁场的坐标表示}绘制为函数图像
\begin{Figure}{载流螺线管的磁场分布图像}
    \inputfigure[scale=0.9]{Chapter10B_Figure09}
\end{Figure}
通过\Refx{图}{载流螺线管的磁场分布图像}我们可以直观的看到通电螺线管内的磁场分布，其中，红线是一个$L,R$比值更大的螺线管，蓝线是一个$L,R$比值更小的螺线管，由此我们可以看出$L,R$比值越大，中央区域越接近匀强磁场，两端磁感应强度的下降曲线越陡。另外也可以看到当$L,R$比值较小时，两端$\mu_0nI/2$的近似程度比中央$\mu_0nI$的近似程度好。

在本节首我们提到过，\textbf{载流螺线管的磁场与条形磁铁是等效的}，下图展现了这点
\begin{Figure}{载流螺线管的磁场与条形磁铁的对比}
    \subfigure[载流螺线管的磁场]
    {
        \includegraphics[scale=0.8]{old/Python/build/VFP_B04.pdf}
    }\qquad
    \subfigure[条形磁铁的磁场]
    {
        \includegraphics[scale=0.8]{old/Python/build/VFP_M04.pdf}
    }
\end{Figure}

\subsubsection{运动电荷的磁场}
在宏观上看，导体中的电荷是以体密度连续分布的，导体中的电流也是连续的，但是如果我们从微观上考察，电荷是若干离散的载流子，电流也就是若干载流子定向运动的结果。

在\Refx{定律}{毕奥-萨伐尔定律}中我们已经了解了电流产生的磁场，电流是载流子定向运动的结果，如果现在我们想通过电流的磁场推出运动电荷的磁场，那么，我们就有必要首先建立起\textbf{宏观连续的电流描述}和\textbf{微观离散的电荷运动描述}之间的关系，这是很有价值的。

第一组关系中，我们将讨论宏观和微观对电荷的不同观点。在宏观上我们认为电荷是连续分布的，可以通过单位体积的电荷量，即电荷体密度$\rho$来描述。在微观上电荷则是一系列离散的带电粒子，显然，带电粒子所带电荷量，带电粒子的分布越集中，其在宏观上就会表现处越大的电荷体密度，记单个电荷的带电量为$q$，记单位体积中的带电粒子数目为$n$，后者也称为\uwave{电荷数密度}或\uwave{电荷浓度}，那么我们就可以将宏观上某一处的电荷体密度$\rho$表示为
\begin{BoxEquation}[电荷的微观表述]
    \rho=nq
\end{BoxEquation}
第二组关系中，我们将讨论宏观和微观对电流的不同观点。为了简单起见，我们假设导体中的电荷和电流分布都是均匀的，若记导体垂直于电流的截面为$S$，那么顺次使用\Refx{公式}{电流与电流密度}和\Refx{定义}{电流密度}以及刚刚的\Refx{公式}{电荷的微观表述}
\begin{equation*}
    I=JS=\rho vS=nqvS
\end{equation*}

这一结果也可以从微观上解释，电流是单位时间内通过导体截面的电荷量，而能在单位时间内通过导体截面$S$的电荷必然分布在\textbf{导体截面前单位时间内电荷的运动距离内}，后者即带电粒子的速度，因而，这一区域的体积是$vS$，这一区域中的带电粒子数目即$nvS$，考虑到每个带电粒子的电荷量为$q$，因而这一区域内的总电荷量即$nqvS$，这也就是电流的大小。
\begin{BoxEquation}[电流的微观表述]
    I=nqvS
\end{BoxEquation}
第三步，我们要建立电流元和电荷运动的关系，我们知道电流元即$I\dif{\vb*{l}}$，其中$\dif{\vb*{l}}$的方向即电流的方向，它的方向等同于$q\vb*{v}$的方向，因此代入\Refx{公式}{电流的微观表述}时
\begin{equation*}
    I\dif{\vb*{l}}=nq\vb*{v}S\dif{l}
\end{equation*}

其中$S\dif{l}$即电流元包含的导体体积，可以记作$\dif{V}$
\begin{equation*}
    I\dif{\vb*{l}}=nq\vb*{v}\dif{V}
\end{equation*}

其中$n$为单位体积的带电粒子数目，故$n\dif{V}$即电流元中的带电粒子数目，记作$\dif{n}$
\begin{BoxEquation}[电流元与运动电荷]
    I\dif{\vb*{l}}=q\vb*{v}\dif{n}
\end{BoxEquation}
这样我们就完整在宏观和微观之间建立起了电流元与运动电荷的关系。现在让我们回到磁场的问题上，若以$\dif{\vb*{B}}$表示电流元的磁场，而以$\vb*{B}$表示电流元中每个运动电荷的磁场
\begin{Equation}[运动电荷的磁感强度1]
    \dif{\vb*{B}}=\vb*{B}nS\dif{l}=\vb*{B}n\dif{V}=\vb*{B}\dif{n}
\end{Equation}
这样就建立了电流元的磁场$\dif{\vb*{B}}$和运动电荷的磁场$\vb*{B}$的关系。

根据\Refx{定律}{毕奥-萨伐尔定律}和\Refx{公式}{电流元与运动电荷}
\begin{Equation}[运动电荷的磁感强度2]
    \dif{\vb*{B}}=
    \frac{\mu_0}{4\pi}
    \frac{I\dif{\vb*{l}}\times\vu*{r}}{r^2}=
    \frac{\mu_0}{4\pi}
    \frac{q\vb*{v}\times\vu*{r}}{r^2}\dif{n}
\end{Equation}
根据式\ref{式:运动电荷的磁感强度1}，即有
\begin{BoxEquation}[运动电荷的磁感强度]
    \vb*{B}=\frac{\mu_0}{4\pi}\frac{q\vb*{v}\times\vu*{r}}{r^2}
\end{BoxEquation}
\Refx{公式}{运动电荷的磁感强度}与\Refx{公式}{点电荷的电场强度}在形式上是对应的，而且也更为简单，那么为何我们最初要选取以电流元为中心的\Refx{定律}{毕奥-萨伐尔定律}作为磁场的起始呢？其中一个原因是，运动电荷构成的恒定电流可以形成静磁场，运动电荷的磁场却不是静磁场，因为作为磁场源的电荷的位置在不断变化，这违背了静磁学的要求。

但是，实际上一个更为重要的原因是，如果现在我们有两个电荷$q_1,q_2$分别以速度$\vb*{v}_1,\vb*{v}_2$在空间中运动，而$\vb*{r}$是$q_2$指向$q_1$的位矢，那么我们可以写出其中$q_1$受到的磁场力和电场力。

电荷$q_1$受到的磁场力$\vb*{F}_B$为
\begin{Equation}[电荷间的电场力]
    \vb*{F}_B=\frac{\mu_0}{4\pi}\frac{q_1q_2\vb*{v}_1\times\vb*{v}_2}{r^2}\times\vu*{r}
\end{Equation}
电荷$q_1$受到的电场力$\vb*{F}_E$为
\begin{Equation}[电荷间的磁场力]
    \vb*{F}_E=\frac{1}{4\pi\varepsilon_0}\frac{q_1q_2}{r^2}\vu*{r}
\end{Equation}
为了便于估计数量级，假定$\vb*{v}_1,\vb*{v}_2$垂直并且$\vb*{v}_2,\vb*{r}$平行，同时
\begin{equation*}
    q_1=q_2=1\si{C}\quad v_1=v_2=1\si{m\cdot s^{-1}}\quad r=1\si{m}
\end{equation*}
代入真空电容率$\varepsilon_0$和真空磁导率$\mu_0$的数值
\begin{equation*}
    \varepsilon_0=8.85\times 10^{-12}\si{C^2\cdot N^{-1}\cdot m^{-2}}\qquad
    \mu_0=4\pi\times 10^{-7}\si{N\cdot A^{-2}}
\end{equation*}
这就可以算得磁场力$F_B=4.0\times 10^{-7}$\si{N}，以及电场力$F_E=9.0\times 10^{9}$\si{N}。

这就是说，两个运动电荷间的电场力要比磁场力大了约$10^{16}$的数量级，因此，通过试探电荷的方式，我们绝无可能分析出另外一个运动电荷产生的磁场，我们会看到运动电荷间的相互作用与它们静止时毫无差别，因为电荷间的电场力要远远大于的磁场力，后者被遮蔽了。

所以\Refx{公式}{运动电荷的磁感强度}虽然很简洁，但其只是一个不能通过实验观测到的理论公式，而电流的磁场是可以实验观测的，因此其无法取代\Refx{定律}{毕奥-萨伐尔定律}作为磁场基本实验定律的地位，就像\Refx{定律}{库仑定律}作为电场基本实验定律那样。

所以为何运动电荷间的磁场力比电场力弱如此之多？以及为何运动电荷间的磁场力无法被观测，但是大量运动电荷构成的电流的磁场就可以被观测？下一小节将解答这两个问题。

\subsection{磁场与电场的统一性}
这一小节让我们来思考一个有趣的问题，我们在式\ref{式:磁场的作用}中曾提及运动电荷受到的磁场力可以表示为$\vb*{F}=q\vb*{v}\times\vb*{B}$，我们可能会疑惑这里的电荷速度$\vb*{v}$是相对于哪一参照系的？

关于选取什么样的参照系作为规定磁场的合适参照系，我们还未曾说过什么。而事实是，任何一个惯性系都可以！对此我们立即会表达出以下的质疑，电荷在不同的惯性参照系中将具有不同的速度，那么按照这里的观点，电荷将受到不同的磁场力，这也就是说电荷受到的磁场力的多少竟然会取决于我们选取的参照系的选取，这听上去实在是有些荒唐。

下面这个例子可以比较清楚的反映这个问题，设有一根电流为$I$在载流长直导线，设导线内的电子以速度$\vb*{v}$向右运动，设导线外有一带负电的粒子同样以速度$\vb*{v}$向右运动。\vspace{5pt}
\begin{Figure}{磁场与电场的统一性}
    \subfigure[基于导线的参照系]
    {
        \inputfigure[scale=0.85]{Chapter10B_Figure11}
    }\\ \vspace{10pt}
    \subfigure[基于电荷的参照系]
    {
        \inputfigure[scale=0.85]{Chapter10B_Figure12}
    }
\end{Figure}
如果我们以导线为参照系，根据$\vb*{F}=qv\times\vb*{B}$可知，此时导线外的电荷将会受到一个指向导线的磁场力，但是，如果我们以电荷为参照系，此时电荷自身的速度变为了零，而静止电荷是不会受到任何磁场力的，然而物体的受力又不应随着参照系的选取而变，因此电荷此时必然将受到一个指向导线的电场力，但是这就意味着运动的“中性”载流导线将带正电。

我们有必要对此仔细检查，为何载流导线在运动时将带电？首先需要肯定的是，电荷量不会随着运动而变化。电荷量的微小差异就会产生巨大的电场力，我们知道物质是中包含带正电的质子和带负电的电子，当我们将物质加热时，由于质子和电子的质量不同，两者因热运动取得的速度增量是不同的，倘若电荷量真与运动有关，两者的电荷量就无法平衡了，而且这种不平衡将会在宏观上可观测，但是我们却从未观察到某种物质在加热后会带电。

我们尚还没有系统的掌握相对论，但解释这里的现象，我们只需要明白以下这样一个简单的事实，\textbf{运动的东西会在运动方向上变得更短}，即\uwave{相对论收缩}，或更为熟知的\uwave{尺缩效应}。

我们知道，在金属导体中，带有负电荷的电子通过运动传递电流，带有正电荷的金属阳离子在材料中是保持不动的，而静止的导线显然是电中性的，即在导线参照系中有$\rho_{+}=\rho_{-}$。

我们现在来考察正负电荷在参照系变换时究竟发生了什么
\begin{itemize}
    \item 以导线为参照系时，正电荷是静止的，以运动电荷为参照系时，正电荷向左运动，根据尺缩原理，新的参照系中正电荷的间距会变小，正电荷的密度$\rho_{+}$会变得更大。
    \item 以导线为参照系时，负电荷向右运动，以运动电荷为参照系时，负电荷是静止的，根据尺缩原理，新的参照系中正电荷的间距会变大，正电荷的密度$\rho_{+}$会变得更小。
\end{itemize}
这就说明在运动电荷的参照系下，这根运动的导线中$\rho_{+}>\rho_{-}$，因而将带有正净电荷，从而对导线外的电荷产生了一个指向导线的电场力，这个力的大小等同于原先的磁场力。

这个例子就告诉我们，电场和磁场是不可分裂或者说是不可分辨的，也就是说，电场和磁场是一个统一整体的两个不同侧面，这个统一整体也就是我们所谓的\uwave{电磁场}。电和磁之间没有严格的界限，电磁场在具体问题中究竟以何种比例表现为电场和磁场，仅取决于我们选取的参照系，就像在这个例子中，在导线参照系中就仅有磁场，在电荷参照系中就仅有电场。

电磁场在我们熟悉的参照系中主要表现为电场，从这样一个观点看，磁场可以被认为是电场由于电荷的运动产生的相对论效应，这种相对论修正的数量级往往在$v^2/c^2$左右，这也是为何先前我们看到运动电荷间的电场力是磁场力的$10^{16}$倍左右，因为光速$c=3\times 10^8$\si{m\cdot s^{-1}}。

磁场是电场的相对论效应，但是，且慢！我们知道相对论效应实际上只有在高速时才会较为显著，在低速时是完全可以忽略的，那么磁场作为一种相对论效应，理论上在我们的低速世界中是不应该被观测到的。然而，实际上我们却可以观测到许多磁效应，例如两根平行载流导线间的磁相互作用力是可以通过实验观测的。磁场作为一种$v^2/c^2$级别的相对论修正对于这里导线中低于$1$\si{m\cdot s^{-1}}的电荷来说本应是可以忽略的，但电流的特殊之处在于，在导线中正负电荷被极为精确的抵消，以至于正常的电场完全消失了，这才使得磁场可以被观测。
\begin{center}
    \small 
    正是电流的这种特殊性质，我们才能在低速世界观测到丰富的磁现象和磁效应。
\end{center}
这也就解释了为何电流的磁场可以被实验观测，而运动电荷的磁场则不行。
